\section{A Taste of Fourier}
\subsection{Fourier Series}

Fourier analysis for a complex function $f$ that is periodic. We let the torus 
$\mathbb{T} = [0, 1) \sim \R / \Z$ and we consider 
\[
    L^2(\mathbb{T}; \C) = \left\{ f : \mathbb{T} \to \C : \int_0^1 |f(x)|^2 dx < \infty \right\} 
    \footnote{Note $|f| = \sqrt{\real(f)^2 + \imag(f)^2}$ }
\] 
Now, $f : \mathbb{T} \to \C$ is the same as $\tilde{f} : \R \to \C$ with $\tilde{f}$
being 1-periodic. We will build an inner product
$(. , .) : L^2(\mathbb{T}; \C) \times L^2(\mathbb{T}; \C) \to \C$.

\begin{defn}[complex inner product]
    A \newword{complex inner product} on a complex vector space $V$ is a function $(. , .) : V \times V \to \C$ 
    such that for all $u, v, w \in V$ and $\alpha \in \C$, we have
    \begin{enumerate}
         \item $(\alpha u + v, w) = \alpha (u, w) + (v, w)$
        \item $(u, v) = \overline{(v, u)}$
        \item $(u, u) \geq 0$ and $(u, u) = 0$ if and only if $u = 0$
    \end{enumerate}
\end{defn}

\begin{prop}
    \[
        (f, g) := \int_0^1 f(x) \overline{g(x)} \; dx
    \] 
    defines a complex inner product on $L^2(\mathbb{T}; \C)$.
\end{prop}
\begin{proof} \leavevmode
    \begin{description}
        \item[(1)] Let $\alpha \in \C$. Then
        \[
            (\alpha f + h, g) = \int_0^1 (\alpha f(x) + h(x)) \overline{g(x)} \; dx = \alpha (f, g) + (h, g)
        \]  
        \item[(2)] 
        \begin{align*}
            (f, g) = \int_0^1 f(x) \overline{g(x)} \; dx &= \int_0^1 (\real f + i \imag f)(\real g - i \imag g) \; dx \\
            &= (a + bi) (c - di) \\
            &= (ac + bd) + i (bc - ad) \\
            &= \overline{(ac + bd) + i (ad - bc)} \\
            &= \overline{(a-bi) (c + di)} \\
            &= \overline{\int_0^1 \overline{(\real f + i \imag f)} (\real g + i \imag g) \; dx} \\
            &= \overline{(g, f)}
        \end{align*}
        \item[(3)] 
        \[
            (f, f) = \int_0^1 f(x) \overline{f(x)} \; dx = \int_0^1 |f(x)|^2 \; dx \geq 0  
        \] 
        Thus $(f, f) = \| f \|_{L^2(\mathbb{T}; \C)}^2 $ and $(f, f) = 0 $ iff $|f| = 0$ for 
        a.e. $x$ iff $f = 0$ in $L^2(\mathbb{T}; \C)$. \qedhere   
    \end{description}
\end{proof}

\begin{rem}
    $f \in L^2(\mathbb{T}; \C) \longleftrightarrow \real (f), \imag (f) \in L^2(\mathbb{T}; \R)$.
    So $\{f_n \}$ is Cauchy in $L^2(\mathbb{T}; \C)$ iff 
$\{ \real (f_n) \}_{n=1}^\infty$ and $\{ \imag (f_n) \}_{n=1}^\infty$ are Cauchy in $L^2(\mathbb{T}; \R)$.

Also, we know $L^2(\mathbb{T}; \R)$ is complete, so $L^2(\mathbb{T}; \C)$ is a Hilbert space.   
\end{rem}
\begin{rem}
    We still have Cauchy Schwartz, i.e. 
    \[
        |(f, g)| \leq \| f \| \| g \|  
    \] 
\end{rem}

\begin{thm}
    Let $e_n(x) := e^{2 \pi i n x} = \cos(2 \pi n x) + i \sin(2 \pi n x)$. Then
    $\{e_{n}(x)\}_{n \in \Z}$ is an orthonormal basis for $L^2(\mathbb{T}; \C)$. 
\end{thm}
\begin{proof}
    Observe 
    \[
        (e_n, e_m) = \int_0^1 e^{2 \pi i n x} e^{-2 \pi i m x} \; dx = \int_0^1 e^{2 \pi i (n - m) x} \; dx
    \] 
    If $n = m$, then $\int_0^1 e^0 \; dx  = 1$, therefore $\| e_n \| = 1 \; \forall n$.

    If $n \neq m$, then 
    \[
    \int_0^1 e^{2 \pi i (n - m) x} \; dx  = \frac{1}{2 \pi i (n - m)} e^{2 \pi i (n - m)x} \Big|_0^1 = \frac{1}{2 \pi i (n - m)} (e^{2 \pi i (n - m)} - 1)
    \]
    But $e^{2 \pi i (n - m)} = \underbrace{\cos(2 \pi (n - m))}_{= 1} + i \underbrace{\sin(2 \pi (n - m))}_{= 0} = 1$.
    
    So $\int_0^1 e^{2 \pi i (n - m) x} \; dx = 0$. Thus $\{e_n\}_{n \in \Z}$ is orthonormal.
    \newline \textbf{Claim:} $\{e_n\}_{n \in \Z}$ is a basis of $L^2(\mathbb{T}; \C)$, i.e. 
    $\forall f \in L^2(\mathbb{T}; \C)$, $\exists \{\alpha_n \}_{n \in \Z} \subseteq \C$ s.t.
    \[
        f = \sum_{n=-\infty}^\infty \alpha_n e_n = \lim_{N \to \infty} \sum_{n=-N}^N \alpha_n e_n
    \]   
    To prove the claim we use complex Stone-Weierstrass. Let
    \[
        \mathcal{A} = \left\{ \sum_{n = -N}^N \alpha_n e_n : \alpha_n \in \C \right\} 
    \] 
    We need that $\mathcal{A}$ is an algebra, which is clear since $e_n e_m = e^{2 \pi i (n + m) x} = e_{n+m}$.
    We also need that $\mathcal{A}$ contains constants, which is clear since $e_0 = 1$. We also
    need that $\mathcal{A}$ is closed under conjugates, which is also clear because $\overline{\alpha_n} \in \C$,
    $\overline{e_n} = e_{-n}$. Finally, we need that $\mathcal{A}$ separates points. 
    If $x \neq y$, $x, y \in [0, 1)$, then $\forall n$ s.t. $nx - ny \notin \Z$, we have
    \[
        e_n(x) = e^{2 \pi i n x} \neq e^{2 \pi i n y} = e_n(y)
    \]
    So $\mathcal{A}$ separates points. 
    
    Thus by complex Stone-Weierstrass, since $\mathbb{T}$ is compact, $\mathcal{A}$ is dense in $C(\mathbb{T}; \C)$
    so $\forall \tilde{f} \in C(\mathbb{T}; \C)$, 
    \[
        \tilde{f} = \lim_{N \to \infty} \sum_{n=-N}^N \alpha_n e_n \quad \text{ limit in }  \| . \|_\infty
    \] 
    Since $\| . \|_\infty \geq \| . \|_{L^2}$, we also have $\tilde{f} = \lim_{N \to \infty} \sum_{n=-N}^N \alpha_n e_n$ in $\| . \|_{L^2}$.
    Finally, we know $C(\mathbb{T}; \C)$ is dense in $L^2(\mathbb{T}; \C)$.
    Therefore $\forall \eta > 0$,
    \[
        \lim_{N \to \infty} \| f - \sum_{n = -N}^N \alpha_n e_n \|_{2} \leq \| f - \tilde{f} \|_2 + 
        \| \tilde{f} - \sum_{n = -N}^N \alpha_n e_n \|_2 \leq \eta + 0 = \eta  
    \]  
    Therefore $f = \sum_{n = -\infty}^\infty \alpha_n e_n$ in $L^2(\mathbb{T}; \C)$, as desired. \qedhere 

\end{proof}

\begin{rem}
    Recall that TFAE:
    \begin{enumerate}
        \item $\{e_n\}_{n \in \Z}$ is a basis 
        \item Every $f \in L^2(\mathbb{T}; \C)$, $f = \sum_{n = -\infty}^\infty \alpha_n e_n = \sum_{n=-\infty}^\infty (f, e_n) e_n$
        \item Completeness: If $(f, e_n) = 0$ for all $n$, then $f = 0$
        \item Parseval: $\| f \|_2^2 = \sum_{n=-\infty}^\infty |(f, e_n)|^2 $  
    \end{enumerate}
\end{rem}

\begin{defn}[Fourier Series]
    For $n \in \Z$, let
    \[
        \hat{f}(n) := (f, e_n) = \int_0^1 f(x) e^{-2 \pi i n x} \; dx
    \]  
    Then the (complex) \newword{Fourier series} of $f : \mathbb{T} \to \C$ is defined by
    \[
        \sum_{n = -\infty}^\infty \hat{f}(n)e^{2 \pi i n x}
    \] 
\end{defn}

\begin{rem}
    It is possible to define a Fourier series for any periodic function. If $f$ is $L-$periodic,
    then 
    \[
        \hat{f}_L(n) := \frac{1}{L} \int_0^L f(x) e^{-2 \pi i n x / L} \; dx
    \]   
    Therefore we get the Fourier series
    \[
        \sum_{n = -\infty}^\infty \hat{f}_L(n) e^{2 \pi i n x / L}
    \] 
\end{rem}
\begin{rem}
    It is also possible to define a real valued Fourier series, using sines and cosines.
    \[
        A_0 + \sum_{n = 1}^\infty A_n \cos(\frac{2 \pi n x}{L}) + B_n \sin(\frac{2 \pi n x}{L})
    \] 
    for some some $A_n, B_n$, also given by inner products.  
\end{rem}

\begin{goal}[question]
    In what sense does 
    \[
        \sum_{n = -\infty}^\infty \hat{f}(n) e^{2 \pi i n x} = \lim_{N \to \infty} \sum_{n = -N}^N \hat{f}(n) e^{2 \pi i n x} =: \lim_{N \to \infty} S_N(x)
    \] 
    converge?
\end{goal}

\begin{prop}
    If $f \in L^2(\mathbb{T}; \C)$, then
    \[
        f(x) = \lim_{N \to \infty} S_N(x) \text{ in } L^2(\mathbb{T}; \C)
    \]  
\end{prop}
\begin{proof}
    Immediate from the basis \qedhere
\end{proof}

Observe that by Parseval, we have 
\[
    \| f \|_2^2 = \sum_{n = -N}^N |(f, e_n)|^2 = \sum_{n=-\infty}^\infty |\hat{f}(n)|^2
\] 
Therefore $\lim_{n \to \infty} |\hat{f}(n)|^2 = 0 \iff \lim_{n \to \infty} |\hat{f}(n)| = 0$. 
\footnote{Can also be seen using Bessel's inequality, $\sum_{n = -\infty}^\infty |\hat{f}(n)|^2 \leq \| f \|_2^2$ }
This gives us

\begin{prop}[Riemann-Lebesque Lemma]
    If $f \in L^2(\mathbb{T}; \C)$, then $\lim_{n \to \infty} |\hat{f}(n)| = 0$.
\end{prop}

\begin{rem}
    This is very useful for the "real Fourier series" version,
    \[
        \lim_{n \to \infty} \int_0^1 f(x) \sin(2 \pi n x) \; dx = \lim_{n \to \infty} \int_0^1 f(x) \cos(2 \pi n x) \; dx = 0
    \] "Oscillatory integrals" (important in PDEs).
\end{rem}

What about other notions of convergence?
\begin{align*}
    S_N(x) &= \sum_{n=-N}^N \hat{f}(n) e^{2 \pi i n x} \\
    &= \sum_{n = -N}^N \left( \int_0^1 f(y) e^{-2 \pi i n y} \; dy \right) e^{2 \pi i n x} \\
    &= \int_0^1 f(y) \left( \sum_{n = -N}^N e^{2 \pi i n (x - y)} \right) dy
\end{align*}
Let $D_N(x) := \sum_{n = -N}^N e^{2 \pi i n x}$. $D_N$ is called the \newword{Dirichlet kernel}.
\[
    S_N(x) = \int_0^1 f(y) D_N(x - y) \; dy = (f * D_N)(x) \quad \text{ "convolution on the torus"}
\] 

\underline{Observe:}
\[
    D_N(x) = 1 + \sum_{n = 1}^N \underbrace{e^{2 \pi i n x}}_{1-periodic} - \underbrace{e^{-2 \pi i n x}}_{\text{1-periodic}}
\]  
Therefore $D_N$ is 1-periodic, and
\[
    \int_0^1 D_N(x) \; dx = \int_0^1 1 + (\text{1-periodic}) \; dx = 1
\]  

\begin{rem}
    $D_N$ is NOT a good kernel. In particular, 
    \[
        \int_0^1 |D_N(x)| \; dx \geq C \log N \to \infty \text{ as } N \to \infty
    \]  
    However, we can rewrite $D_N$ as
    \begin{align*}
        D_N(x) = \sum_{n = -N}^N e^{2 \pi i n x}  &= \sum_{n = -N}^N \left(e^{2 \pi i x}\right)^n\\
        &= \sum_{n = 0}^{2N} e^{2 \pi i x (n - N)} \\
        &= e^{-2 \pi i N x} \underbrace{\sum_{n = 0}^{2N} \left(e^{2 \pi i x}\right)^n}_{\text{finite geometric series}} \\
        &= e^{-2 \pi i N x} \frac{1 - e^{2 \pi i x (2N + 1)}}{1 - e^{2 \pi i x}} \\
        &= \frac{e^{-2 \pi i N x} - e^{2 \pi i x (N + 1)}}{1 - e^{2 \pi i x}} \\
        &= \frac{e^{-2 \pi i x \cdot \frac{1}{2}} \left( e^{-2 \pi i N x } - e^{2 \pi i (N + 1) x } \right)}{e^{-2 \pi i x \cdot \frac{1}{2}}\left(1 - e^{2 \pi i x}\right)} \\
        &= \frac{e^{-2 \pi i (N + \frac{1}{2}) x} - e^{-2 \pi i (N + \frac{1}{2}) x}}{e^{-2 \pi i \frac{1}{2} x} - e^{2 \pi i \frac{1}{2} x}} \\
        &= \frac{\sin(2 \pi (N + \frac{1}{2} )x)}{\sin(\pi x)}
    \end{align*}
    So $D_N$ is a "rescaled sine function". The sine function is a good kernel, but the rescaling is bad.
    Using this representation, we have the following theorem.
\end{rem}

\begin{thm}[pointwise convergence]
    Let $f \in L^2(\mathbb{T}; \C)$, and $f$ is Lipschitz continuous at $x_0$. Then
    \[
        S_N f(x_0) \to f(x_0) \text{ as } N \to \infty
    \]   
\end{thm}
\begin{proof}
    Proof on the assignment
\end{proof}
\begin{thm}[uniform convergence]
    If $f \in C^2(\mathbb{T}; \C)$ then $S_N f (x) \to f(x)$  on $\mathbb{T}$.   
\end{thm}
\begin{proof}
    Proof on the assignment.
\end{proof}

\begin{rem}
    In fact, we see
    \[
        \hat{f}(n) = \int_0^1 f(x) e^{-2 \pi i n x} \; dx
    \]
    is well defined whenever $f \in L^1(\mathbb{T}; \C)$. \footnote{Since $|e^{-2 \pi i n x}| = 1$}
    Therefore many Fourier things will be good for $L^1(\mathbb{T}; \C)$, not just $L^2(\mathbb{T}; \C)$. 
\end{rem}

\subsection{The Fourier Transform}
Last class, we had $\hat{f} = (f, e_n) = \int_0^1 f(x) e^{-2 \pi i n x} \; dx$ 
are the fourier coefficients for $f : \mathbb{T} \to \C$, and this quanity is finite $\forall
f \in L^1(\mathbb{T}; \C)$, $L^2(\mathbb{T}; \C) \subseteq L^1(\mathbb{T}; \C)$. We now introduce:

\begin{defn}[Fourier Transform]
    Let $f : \R \to \C$. Then $\forall \zeta \in \R$,
    \[
        \hat{f}(\zeta) := \int_{\R} f(x) e^{-2 \pi i \zeta x} \; dx
    \]   
\end{defn}

\begin{rem}
    We will reserve $x, y, z$ for variables in the physical space, and $\zeta, \eta, \xi$ for
    variables in the "Fourier space".  
\end{rem}

\begin{prop}
    If $f \in L^1(\R; \C)$, then $\hat{f} \in L^\infty(\R; \C) \cap C(\R; \C)$. 
    \footnote{We will drop the $\C$ from now on, but everything can be complex valued. }
\end{prop}
\begin{proof}
    We have
    \begin{align*}
        |\hat{f}(\xi)| = \left| \int_{\R} f(x) e^{-2 \pi i \xi x} \; dx \right| \leq 
        \int_{\R} |f(x)| |e^{-2 \pi i \xi x}| \; dx = \int_{\R} |f(x)| \; dx = \| f \|_1
    \end{align*}
    Therefore $\| \hat{f} \|_\infty \leq \| f \|_1 $.

    Moreover, $\forall h > 0$,
    \begin{align*}
        |\hat{f}(\zeta + h) - \hat{f}(\zeta)| &= \left| \int_{\R} f(x) e^{-2 \pi i (\zeta + h) x} \; dx - \int_{\R} f(x) e^{-2 \pi i \zeta x} \; dx \right| \\
        &= \left| \int_{\R} f(x) e^{-2 \pi i \zeta x}(e^{-2 \pi i h x} - 1) \; dx \right| \\
        &\leq \int_{\R} |f(x)| \underbrace{|e^{-2 \pi i \zeta x}|}_{=1} |e^{-2 \pi i h x} - 1| \; dx
    \end{align*}
    Also, $\lim_{h \to 0} |e^{-2 \pi i h x} - 1| = 0$, $\forall x \in \R$ and
    \[
        |f(x)| |e^{-2 \pi i h x} - 1| \leq |f(x)|(|e^{-2 \pi i h x}| + 1) = 2 |f(x)| \in L^1(\R)
    \]  
    Therefore by DCT,
    \[
        \lim_{h \to 0} |\hat{f}(\zeta + h) - \hat{f}(\zeta)| = \int_{\R} |f(x)| \underbrace{\lim_{h \to 0} |e^{-2 \pi i h x} - 1|}_{=0} \; dx = 0 \qedhere
    \] 
\end{proof}

\begin{prop}
    Let $f,g \in L^1(\R)$. Then
    \begin{enumerate}
        \item $\widehat{f*g} = \hat{f} \cdot \hat{g}$;
        \item $\int_{\R} f \hat{g} \; dx = \int_{\R} \hat{f} g \; dx$ 
    \end{enumerate} 
\end{prop}
\begin{proof} \leavevmode
    \begin{description}
        \item[(a)] 
        \[
            \widehat{f * g}(\zeta) = \int (f*g)(x) e^{-2 \pi i \zeta x} \; dx = \iint f(x-y) g(y) \; dy \; e^{-2 \pi i \zeta x} \; dx
        \]  
        Since $f, g \in L^1(\R)$, this integral is finite\footnote{b/c $\leq \| f \|_1 \| g \|_1$ } so 
        by Fubini,
        \[
            = \iint_{\R} f(x-y) e^{-2 \pi i \zeta (x-y)} \; dx \; g(y) e^{-2 \pi i \zeta y} \; dy
        \] 
        Using $z = x-y$, 
        \[
            = \int \left( \int_{\R} f(z) e^{-2 \pi i \zeta z} \; dz \right) g(y) e^{-2 \pi i \zeta y} \; dy
            = \hat{f}(\zeta) \int g(y) e^{-2 \pi i \zeta y} \; dy = \hat{f} \hat{g}
        \]  

        \item[(b)] 
        \begin{align*}
            \int f(x) \hat{g(x)} &= \int f(x) \int g(y) e^{-2 \pi i x y} \; dy \; dx \\
            &= \int g(y) \underbrace{\int f(x) e^{-2 \pi i x y} \; dx}_{\hat{f}(y)} \; dy \quad \text{by Fubini} \\ 
            &= \int \hat{f}(y) g(y) \; dy \qedhere \\
        \end{align*}
    \end{description}
\end{proof}

\begin{lem}
    Let $f(x) = e^{-\pi a x^2}$, $a > 0$. Then
    \[
        \hat{f}(\zeta) = \frac{1}{\sqrt{a}}e^{-\pi \zeta^2 / a}
    \]   
\end{lem}
\begin{proof}
    We compute 
    \begin{align*}
        \hat{f'}(\zeta) &= \int_{\R} f'(x) e^{-2 \pi i \zeta x} \; dx \\
        &= f(x) e^{-2 \pi i \zeta x} \Big|_{-\infty}^\infty - \int_{\R} f(x) (-2 \pi i \zeta) e^{-2 \pi i \zeta x} \; dx \quad \text{by IBP}\\
        &= e^{-\pi a x^2}e^{-2 \pi i \zeta x} \Big|_{-\infty}^\infty + 2 \pi i \zeta \int_{\R} f(x) e^{-2 \pi i \zeta x} \; dx \\
        &= 0 - 0 + 2 \pi i \zeta \hat{f}(\zeta)
    \end{align*}
    Therefore $\hat{f'}(\zeta) = 2 \pi i \zeta \hat{f}(\zeta)$
    
    Also, 
    \[
        \frac{d}{d\zeta} \hat{f}(\zeta) = \frac{d}{d \zeta} \left( \int_{\R} f(x) e^{-2 \pi i \zeta x} \; dx \right)
    \] 
    To differentiate under the integral, we need to check DCT.
    \[
        \left| f(x) \left( \frac{e^{-2 \pi i x (\zeta + h)} - e^{-2 \pi i x \zeta}}{h} \right) \right| \leq |f(x)| \left| \frac{d}{d \zeta} e^{-2 \pi i x \zeta} \right| \Big|_{\zeta = \zeta^*} 
    \] 
    where $\zeta^* \in [\zeta, \zeta + h]$. Then
    \begin{align*}
        = |f(x)| |-2 \pi i x| |e^{-2 \pi i x \zeta^*}| = 2 \pi |f(x)| |x| = 2 \pi |x| e^{-\pi a x^2} \in L^1(\R)
    \end{align*}
    By DCT, 
    \begin{align*}
        \frac{d}{d \zeta} \hat{f}(\zeta) &= \int_{\R} f(x) \frac{d}{d \zeta}(e^{-2 \pi i x \zeta}) \; dx \\
        &= \int_{\R} f(x) (-2 \pi i x) e^{-2 \pi i x \zeta} \; dx \\
        &= i \int_{\R} (-2 \pi x) e^{-\pi a x^2} e^{-2 \pi i x \zeta} \; dx \\
        &= \frac{i}{a} \int_{\R} (e^{-\pi a x^2})' e^{-2 \pi i x \zeta} \; dx \\
        &= \frac{i}{a} \underbrace{\int_{\R} f'(x) e^{-2 \pi i \zeta x} \; dx}_{\hat{f'}(\zeta)} \\
        &= \frac{1}{a} 2 \pi i \zeta \hat{f}(\zeta) \\
        &= -\frac{2 \pi \zeta}{a} \hat{f}(\zeta)
    \end{align*}
    Thus,
    \begin{align*}
        \frac{d}{d \zeta} (e^{\pi \zeta^2 / a} \hat{f}(\zeta)) &= e^{\pi \zeta^2 / a} \frac{d}{d \zeta} \hat{f}(\zeta) + \frac{d}{d \zeta}(e^{\pi \zeta^2 / a}) \hat{f}(\zeta) = 0 \\
        &= \frac{-2 \pi \zeta}{a} e^{\pi \zeta^2 / a} \hat{f}(\zeta) + \frac{2 \pi \zeta}{a} e^{\pi \zeta^2/a} \hat{f}(\zeta) = 0
    \end{align*}
    So $\zeta \mapsto e^{\pi \zeta^2 / a} \hat{f}(\zeta)$ is constant in $\zeta$. Therefore
    \[
        e^{\pi \zeta^2/a} \hat{f}(\zeta) = e^{0} \hat{f}(0) = \hat{f}(0) = \int_{\R} e^{-\pi a x^2} \; dx = \frac{1}{\sqrt{a}} \int_{\R} e^{-\pi y^2} \; dy = \frac{1}{\sqrt{a}}
    \]  
    Therefore $\hat{f}(\zeta) = \frac{1}{\sqrt{a}} e^{-\pi \zeta^2 / a}$. \qedhere 
\end{proof}

\begin{defn}[Inverse Fourier Transform]
    If $f \in L^1(\R)$, then 
    \[
        \check{f}(x) := \int_{\R} f(\zeta) e^{2 \pi i x \zeta} d \zeta = \hat{f}(-x)
    \] 
\end{defn}
\begin{proof}
    If $f \in L^1(\R)$, then $\check{f} \in L^\infty(\R) \cap C(\R)$. The proof is identical to $\hat{f}$.   
\end{proof}

\begin{rem}
    We hope that if $f \in L^1(\R)$, then $\check{\hat{f}} = f$. However,
    \[
        \check{\hat{f}} = \int \left( \int f(z) e^{-2 \pi i \zeta z} \; dz \right) e^{2 \pi i x \zeta} \; d\zeta
        \leq \iint |f(z)| \; dz \; d \zeta = \infty
    \] 
\end{rem}
\begin{thm}[Fourier Inversion]
    If $f \in L^1(\R)$ and $\hat{f} \in L^1(\R)$, then $\exists f_0 \in C(\R)$ s.t. $f = f_0$ almost everywhere, and
    $f_0(x) = \hat{\check{f}} = \check{\hat{f}}.$ 
\end{thm}

\begin{rem}
    The inversion formula is an analog of convergence of Fourier series, because it says
    \[
        f(x) = \int \hat{f}(\zeta) e^{2 \pi i \zeta x} \; d\zeta \approx \sum_{\zeta \in \R} \hat{f}(\zeta) e^{2 \pi i \zeta x}
    \]  

\end{rem}

\begin{cor}
    If $ f \in L^1(\R)$, $\hat{f} \equiv 0 \Rightarrow f \equiv 0$ a.e. 
\end{cor}
\begin{proof}[Proof of Thm]
    Let $\epsilon > 0$. Fix $x \in \R$, and define $\phi^\epsilon(\zeta) = e^{2 \pi i x \zeta} e^{- \pi \epsilon \zeta^2}$ 
    Then
    \begin{align*}
        \hat{\phi^{\epsilon}}(y) &= \int \phi^{\epsilon}(\zeta) e^{-2 \pi i y \zeta} \; d \zeta \\
        &= \int e^{2 \pi i x \zeta} e^{- \pi \epsilon \zeta^2} e^{-2 \pi i y \zeta} \; d \zeta \\
        &= \int e^{- \pi \epsilon \zeta^2} e^{- 2 \pi i (y - x) \zeta} \; d \zeta \\
        &= \widehat{e^{-\pi \epsilon (.)^2}}(y - x) \\
        &= \frac{1}{\sqrt{\epsilon}} e^{-\pi (y - x)^2 / \epsilon} \quad \text{by lemma}\\
    \end{align*}   
    So 
    \[
        \int f \hat{\phi^\epsilon} \; dy = \int \hat{f} \phi^\epsilon \; dy
    \] 
    Therefore 
    \[
        \int f(y) \frac{1}{\sqrt{\epsilon}} e^{-\pi (x-y)^2 / \epsilon} \; dy = \int \hat{f}(y) \phi^\epsilon(y) \; dy
    \] 
    Let $K_\epsilon(y) := \frac{1}{\sqrt{\epsilon}} e^{-\pi y^2 / \epsilon}$. We saw that $\{ K_\epsilon\}$ are good kernels,
    so 
    \[
        \int \hat{f}(y) \phi^{\epsilon}(y) \; dy = (K_\epsilon * f)(x)
    \]   
    \textbf{Claim:} $K_\epsilon * f \to f$ in $L^1(\R)$.
    \newline \noindent If the claim holds, then
    \[
        f(x) = \lim_{\epsilon \to 0} \int \underbrace{\hat{f}(y) e^{2 \pi i x y} e^{- \pi \epsilon y^2}}_{\leq \|\hat{f}\|_{L^1}} \; dy
    \]  
    So by DCT,
    \[
        f(x) = \int \hat{f}(y) e^{2 \pi i x y} \; dy 
    \]  
    $f(x) = \check{\hat{f}}(x)$, therefore $f = \check{\hat{f}}$ a.e.
    An identical proof gives $f = \hat{\check{f}}$ a.e. Moreover, $\check{\hat{f}}$ and $\hat{\check{f}}$ are continuous,
    so $\exists f_0 \in C(\R)$ s.t. $f_0 = f$ a.e.
    \newline \noindent \textbf{Proof of claim:} 
    \newline \noindent \begin{align*}
        |(f * K_\epsilon)(x) - f(x)| &= \left| \int (f(x-y) - f(x)) K_\epsilon(y) \; dy \right| \\
        &\leq \int |f(x-y) - f(x)| K_\epsilon(y) \; dy \\ 
    \end{align*}
    Therefore
    \begin{align*}
        \int |(f * K_\epsilon)(x) - f(x)| \; dx &\leq \iint |f(x-y) - f(x)| K_\epsilon(y) \; dy \; dx \\
        &= \iint |f(x-y) - f(x) | \frac{1}{\sqrt{\epsilon}} e^{-\pi (y)^2 / \epsilon} \; dy \; dx \\
        &= \iint |f(x - \sqrt{\epsilon}y) - f(x)| \; dx \; e^{-\pi y^2} \; dy
    \end{align*} 
    Then by two DCT's, we can take $\epsilon \to 0$ inside the integral. (Finish proof in the posted notes)
\end{proof}

Let us introduce
\[
    C_0(\R) = \left\{ f \in C(\R) : \lim_{|x| \to \infty} |f(x)| = 0 \right\} \footnote{So $C_c(\R) \subseteq C_0(\R)$ }
\] 

\begin{prop}
    If $f \in C^1(\R) \cap C_0(\R) \cap L^1(\R)$ and $f' \in L^1(\R)$, then
    \[
        \widehat{f'}(\zeta) = (2 \pi i \zeta) \hat{f}(\zeta)
    \] 
    Additionally, if $f \in L^1(\R)$, then $\hat{f} \in C_0(\R)$. 
\end{prop}

\begin{proof}
    Identical to before. 
    \begin{align*}
        \widehat{f'}(\zeta) &= \int_{\R} f'(x) e^{-2 \pi i \zeta x} \; dx \\
        &= f(x) e^{-2 \pi i \zeta x} \Big|+{x = -\infty}^\infty - \int_{\R} f(x) (-2 \pi i \zeta) e^{-2 \pi i \zeta x} \; dx \\
        &= 0 - 0 + 2 \pi i \zeta \hat{f}(\zeta)
    \end{align*}
    For the second part, if $f \in L^1(\R)$, then $\exists \{f_n\} \subseteq C_c^\infty(\R)$ s.t. $\| f_n - f \|_1 \to 0$.
    Then $f_n' \in C_c(\R) \; \forall n$, therefore $f_n' \in L^1(\R)$. Therefore $\hat{f_n'} \in L^\infty(\R) \cap C(\R) \; \forall n$.
    So by the first claim,
    \[
        \hat{f_n'}(\zeta) = 2 \pi i \zeta \hat{f_n}(\zeta) \in L^\infty(\R) \cap C(\R) \; \forall n
    \]    
    Since $2 \pi i (.) \hat{f_n}(.) \in L^\infty(\R)$, we have
    \[
        |2 \pi i \zeta \hat{f_n}(\zeta)| \leq C, \Rightarrow |\zeta| |\hat{f_n}(\zeta)| \leq C \Rightarrow |\hat{f_n}(\zeta)| \leq \frac{C}{|\zeta|} \to 0 \text{ as } |\zeta| \to \infty
    \]  
    Therefore $\hat{f_n} \in C_0(\R), \; \forall n$.
    
    Moreover, since $\| f_n - f \|_1 \to 0 $, 
    \begin{align*}
        \| \hat{f_n} - \hat{f} \|_\infty &= \sup_{\zeta} \left| \int f_n(x) e^{-2 \pi i \zeta x} \; dx - \int f(x) e^{-2 \pi i \zeta x} \; dx \right| \\
        &= \sup_{\zeta} \left| \int (f_n(x) - f(x)) e^{-2 \pi i \zeta x} \; dx \right| \\
        &\leq \| f_n - f \|_1 \xrightarrow{n \to \infty} 0    
    \end{align*} 

    Finally, $\forall \zeta \in \R$,
    \[
        |\hat{f}(\zeta)| \leq |\hat{f_n}(\zeta)| + \| \hat{f_n} - \hat{f} \|_\infty 
    \]  
    Therefore
    \[
        \lim_{|\zeta| \to \infty} |\hat{f}(\zeta)| \leq \underbrace{\lim_{|\zeta| \to \infty} |\hat{f_n}(\zeta)|}_{0 \text{ bc } \hat{f_n} \in C_0(\R)} + \| \hat{f_n} - \hat{f} \|_\infty 
    \] 
    Taking $n \to \infty$,
    \[
        \lim_{|\zeta| \to \infty} |\hat{f}(\zeta)| = 0 \quad \Rightarrow \quad \hat{f} \in C_0(\R) \qedhere
    \]  
\end{proof}

\begin{rem}
    The second part of the proposition is the analogue of the Riemann-Lebesgue lemma except $f \in L^1(\R)$.
    \[
        \left[\lim_{n \to \infty} \int_0^1 f(x) e^{-2 \pi i n x} \; dx = 0 \right]
    \]      
\end{rem}

This gives us a good theory of Fourier Transforms for $f \in L^1(\R)$. However we would really like
to do it for $f \in L^2(\R)$ like before, and there's no easy way to compare $L^2(\R)$ and $L^1(\R)$.

\begin{thm}[Plancherel's Theorem]
    If $f \in L^1(\R) \cap L^2(\R)$, then $\hat{f} \in L^2(\R)$ and
    \[
        \| f \|_2 = \| \hat{f} \|_2 \quad \text{(isometry)}
    \]   
\end{thm}
\begin{rem}
    Recall Parseval's identity for Fourier series. 
    \[
        \| f \|_2^2 - \sum_{n \in \Z} |\hat{f_n}(n)|^2 \quad (\ell^2 \text{ version of Plancherel})
    \] 
\end{rem}
\begin{proof}
    Let $f \in L^1(\R) \cap L^2(\R)$. Let $g(x) := \overline{f(-x)}$. Then $g \in L^1(\R) \cap L^2(\R)$.
    Consider $\omega(x) := (f*g)(x)$. Then by Young's inequality,
    \[
        \| \omega \|_1 \leq \| f \|_1 \| g \|_1 \quad \Rightarrow \omega \in L^1(\R)
    \] 
    \textbf{Claim:} $\omega$ is continuous at $0$.
    \newline \noindent \begin{align*}
        |\omega(h) - \omega(0)| &= \left| \int f(h-y)g(y) \; dy - \int f(0 -y) g(y) \; dy \right| \\
        &= \left| \int (f(h-y) - f(-y))g(y) \; dy  \right| \\
        &\leq \| f(h - .) - f(.) \|_2 \| g \|_2 \\
        &= \| f(. - h) - f(.) \|_2 \| g \|_2 = \| \tau_h f - f \|_2 \| g \|_2        
    \end{align*}  
    Since $f \in L^2(\R)$, $\forall \eta > 0$, $\exists \tilde{f} \in C_c(\R)$ s.t. $\| \tilde{f} - f \|_2 < \eta$.
    Then,
    \[
        \| \tau_h f - f \|_2 \leq \underbrace{\| \tau_h f - \tau_h \tilde{f} \|_2}_{= \| f - \tilde{f} \|_2 } + \| \tau_h \tilde{f} - \tilde{f} \|_2 + \| \tilde{f} - f \|_2    
        \leq 2 \eta + \underbrace{\| \tau_h \tilde{f} - \tilde{f} \|_2}_{\tilde{f} \in C_c(\R)}
    \]    
    Therefore
    \[
        \| \tau_h \tilde{f} - \tilde{f} \|_\infty \xrightarrow{h \to 0} 0 \quad \Rightarrow \quad \| \tau_h \tilde{f} - \tilde{f} \|_2 \xrightarrow{h \to 0} 0 
    \] 
    So 
    \[
        \lim_{h \to 0} \| \tau_h f  - f \|_2 \leq 2 \eta \quad \forall \eta > 0 \quad \Rightarrow \quad \lim_{h \to 0} \| \tau_h f - f \|_2 = 0 
    \]  
    Thus $\lim_{h \to 0} |\omega(h) - \omega(0)| = 0$, and thus $\omega$ is continuous at $0$.

    Since $ \omega = f *g$, we have $\hat{\omega} = \hat{f}\hat{g}$ and 
    \begin{align*}
        \hat{g}(\zeta) = \widehat{f(- \cdot)}(\zeta) &= \int_{\R} \overline{f(-x)}e^{-2 \pi i \zeta x} \; dx \\
        &= \int_{\R} \overline{f(-x)e^{2 \pi i \zeta x}} \; dx \\
        &= \overline{\int{\R} f(-x) e^{2 \pi i \zeta x} \; dx} \\
        &= \overline{\int_{\R} f(x) e^{-2 \pi i \zeta x} \; dx} \\
        &= \overline{\hat{f}(\zeta)}
    \end{align*}
    So $\hat{\omega} = \hat{f} \hat{g} = \hat{f} \overline{\hat{f}} = |\hat{f}|^2 \geq 0$
    Finally, recall the good kernel
    \[
        K_\epsilon(y) = \frac{1}{\sqrt{\epsilon}} e^{-\pi y^2 / \epsilon} = \widehat{e^{-\pi \epsilon (.)^2}}(y)
    \]  
    Using $\int \hat{\omega} f \; dy = \int \omega \hat{f} \; dy$, we get
    \begin{align*}
        \int \hat{\omega}(y) e^{-\pi \epsilon y^2} \; dy &= \int \omega(y) \frac{1}{\sqrt{\epsilon}e^{-\pi y^2 / \epsilon}} \; dy \\
        &= \int \omega(y) K_\epsilon(y) \; dy \\
        &= \int \omega(y) K_\epsilon(-y) \; dy \\
        &= (\omega * K_\epsilon)(0) \quad (*)
    \end{align*} 
    On the LHS, $\hat{\omega} \geq 0$, and $\hat{\omega}(y)e^{-\pi \epsilon y^2} \to \hat{\omega}(y) $ as $\epsilon \to 0$.
    So by MCT, 
    \[
        \lim_{\epsilon \to 0} \int \hat{\omega}(y) e^{- \pi \epsilon y^2} \; dy = \int \hat{\omega}(y) \; dy
    \] 
    \textbf{Claim:} $\lim_{\epsilon \to 0} (\omega * K_\epsilon)(0) = \omega(0)$
    \newline \noindent If claim holds, then by $(*)$, 
    \begin{align*}
        \int \hat{\omega}(y) \; dy = \omega(0) \quad \Rightarrow \quad \int |\hat{f}(y)|^2 \; dy &= \int f(y) \overline{f(-(0 - y))} \; dy \\
        &= \int f(y) \overline{f(y)} \; dy \\
        &= \int |f(y)|^2 \; dy
    \end{align*}  
    So $\| f \|_2^2 - \| \hat{f} \|_2^2  $ and we are done.
    
    Now to prove the claim. Fix $ \eta > 0$. Since $\omega$ is continuous at $0$, $\exists \delta > 0$ s.t.
    if $|y| < \delta$, then $|\omega(y) - \omega(0)| < \eta$. Therefore
    \begin{align*}
        &\left| \int \omega(0 - y) K_\epsilon(y) \; dy - \omega(0) \right| = \left| \int (\omega(0 -y) - \omega(0))K_\epsilon(y) \; dy \right|  \\
        &\leq \int_{|y| \leq \delta} |\omega(-y) - \omega(0)| K_\epsilon(y) \; dy  + \int_{|y| > \delta} (|\omega(-y)| _ |\omega(0)|) K_\epsilon(y) \; dy \marginnote{$|\omega(0)| < \infty$ because $\omega$ is continuous at $0$}\\
        &\leq \underbrace{\int_{|y| \leq \delta} \eta K_\epsilon(y) \; dy}_{\leq \eta} + \underbrace{|\omega(0)|\int_{|y| > \delta} K_\epsilon(y) \; dy}_{\to 0 \text{ as } \epsilon \to 0}  + \int_{|y| > \delta} |\omega(-y)| K_\epsilon(y) \; dy \\
    \end{align*}      
    Finally, 
    \begin{align*}
        \int_{|y| > \delta} |\omega(-y)| \frac{1}{\sqrt{\epsilon}}e^{- \pi y^2 / \epsilon} \; dy &\leq \int_{|y| > \delta} |\omega(-y)| \frac{1}{\sqrt{\epsilon}} e^{-\pi \delta^2 / \epsilon} \; dy \\
        &\leq \underbrace{\frac{1}{\sqrt{\epsilon}}e^{- \pi \delta^2 / \epsilon}}_{\to 0 \text{ as } \epsilon \to 0} \underbrace{\int |\omega(-y)| \; dy}_{= \| \omega \|_1 } \\
    \end{align*}
    So we've shown $\forall \eta > 0$,
    \[
        \lim_{\epsilon \to 0} | (\omega * K_\epsilon)(0) - \omega(0)| \leq \eta
    \]  
    Therefore $\lim_{\epsilon \to 0} | (\omega * K_\epsilon)(0) - \omega(0)| = 0$, and we are done. \qedhere
\end{proof}

\begin{defn}[Fourier Transform on $L^2(\R)$]
    Let $f \in L^2(\R)$. Then $\exists \{ f_k\} \subseteq C_c^\infty(\R)$ s.t. $f_k \to f$ in $L^2(\R)$.
    Since $\{f_k\} \subseteq C_c^\infty(\R)$, $f_k \in L^2(\R) \cap L^1(\R), \forall k$. \footnote{Smooth and compact functions are in every single $L^p$ space. }      
    Since $\{f_k\}$ converge in $L^2(\R)$, they are also Cauchy in $L^2(\R)$, and
    \[
        \underbrace{\| f_j - f_k \|_2}_{\in L^2(\R) \cap L^1(\R)} =\footnote{By the Plancherel theorem} \| \widehat{f_j - f_k} \|_2
        = \| \hat{f_j} - \hat{f_k} \|_2 \footnote{By linearity of Fourier transform} \quad \forall j, k
    \]    
    Therefore $\{\hat{f_j}\}$ is Cauchy in $L^2(\R)$, so by the completeness of $L^2(\R)$, 
    we define
    \[
        \hat{f}(\zeta) := \lim_{j \to \infty} \hat{f_j}(\zeta) \text{ in } L^2(\R)
    \]    
    \textbf{Claim:} $\hat{f}$ is well-defined.
    \newline \noindent Suppose $\exists \{f_k\}, \{f_k^*\} $ s.t. $f_k \to f$ and $f_k^* \to f$ in $L^2(\R)$.
    Then we claim $\| \hat{f} - \hat{f^*} \|_2 = 0$. We have
    \begin{align*}
        \| f_k - f_k^* \|_2 &\leq \| f_k - f \|_2 + \| f - f_k^* \|_2 \xrightarrow{k \to \infty} 0 \\  
    \end{align*}   
    And by Plancherel
    \[
        \| f_k - f_k^* \|_2 = \| \hat{f_k} - \hat{f_k^*} \|_2
    \] 
    Therefore $\| \hat{f_k} - \hat{f_k^*} \|_2 \to 0$ as $k \to \infty$. So $\hat{f} = \hat{f^*}$ in $L^2(\R)$, and thus
    $\hat{f}$ is well-defined.
\end{defn}
\begin{prop}
        Let $f \in L^2(\R) \cap C^1(\R)$. If $f' \in L^2(\R)$, and $(\cdot) \hat{f}(\cdot) \in L^2(\R)$,
        then $\widehat{f'}(\zeta) = 2 \pi i \zeta \hat{f}(\zeta)$. \footnote{Before, we had the same result but with $f \in L^1(\R)$ and $f \in C_0(\R)$.}   
\end{prop}
\begin{proof}
    Let $g_n(x) := (p_{1/n} * f)(x)$. Then $g_n \in C^\infty(\R)$, and $g_n \to f$ in $L^2(\R)$.
    Since $f \in C^1(\R)$, we have $f' \in C(\R)$. Then
    \[
        g_n'(x) = \frac{d}{dx} \left( \int p_{1/n}(y)f(x-y) \; dy \right)
    \] 
    By DCT, 
    \begin{align*}
        \left| p_{1/n}(y) \frac{f(x + h - y) - f(x-y)}{h}\right| &\leq p_{1/n}(y) |f'(x^*)| \quad \text{for some } x^* \in (x-h, x+h) \footnote{By the mean value theorem}\\
        &\leq \underbrace{\sup_{B(x^*, 1)} |f'| p_{1/n}(y)}_{\in L^1(\R)} \\
    \end{align*}  
    Thus $\forall x$,
    \[
        g_n'(x) = \int p_{1/n}(y)f'(x-y) \; dy
    \]  
    and since $f' \in L^2(\R)$, $\lim_{n \to \infty} \| g_n' - f' \|_2 = 0$.
    So $g_n$ are not in $C_0(\R)$, so we do one more level of approximation. 
    
    Let $\gamma \in C_c^\infty(\R)$ s.t. 
    \[
        \gamma(x) = \begin{cases}
        1 & \text{ if } |x| \leq 1 \\
        0 & \text{ if } |x| \geq 2
        \end{cases}
    \]  
    Let $\gamma_n(x) := \gamma(\frac{x}{n})$. Then $|\gamma_n| \leq 1$, $|\gamma_n'| \leq \frac{1}{n}|\gamma'| \leq \frac{c}{n}$.
    \textbf{Claim:} $\| \gamma_n f - f \|_2 \to 0 $.
    \newline \noindent 
    \[
        \int \left| \gamma_n f(x) - f(x) \right|^2 \; dx
    \]
    $\gamma_n f \to f$ pointwise, and $\left| \gamma_n f(x) - f(x) \right|^2 \leq 2 |\gamma_n f|^2 + 2|f|^2 \leq 4|f|^2 \in L^1(\R)$.
    So by DCT, $\lim_{n \to \infty} \| \gamma_n f - f \|_2 = 0$ 
    
    Finally, let $f_n := \gamma_n g_n$. Then $f_n \in C_c^\infty(\R)$, and
    \begin{align*}
        \| f_n - f \|_2 = \| \gamma_n g_n - f \|_2 &\leq \| \gamma_n g_n - \gamma_n f \|_2 + \| \gamma_n f - f \|_2 \\
        &= \| \gamma_n (g_n - f) \|_2 + \| \gamma_n f - f \|_2 \\
        &\leq \| g_n - f \|_2 + \| \gamma_n f - f \|_2 \xrightarrow{n \to \infty} 0        
    \end{align*}   
    Finally, $f_n'(x) = \gamma_n(x) g_n'(x) + \gamma_n'(x) g_n(x)$.
    \begin{align*}
        \| f_n' - f' \|_2 &= \| \gamma_n g_n' + \gamma_n' g_n - f' \|_2 \\
        &\leq \| \gamma_n' g_n \|_2 + \| \gamma_n g_n' - \gamma_n f' \|_2 + \| \gamma_n f' - f' \|_2 \\
        &\leq \| \gamma_n' \|_\infty \| p_{1/n} * f \|_2 + \| \gamma_n \|_\infty \| g_n' - f' \|_2 + \| \gamma_n f' - f' \|_2 \\
        &\leq \frac{c}{n} \| f \|_2 \| p_{1/n} \|_1 + \| g_n' - f' \|_2 + \| \gamma_n f' - f' \|_2 \xrightarrow{n \to \infty} 0            
    \end{align*} 
    So we have $\{f_n\} \subseteq C_c^\infty(\R)$ s.t. $f_n \to f$ and $f_n' \to f'$ in $L^2(\R)$.
    Since $f_n \in C_c^\infty(\R)$, $f_n \in L^1(\R) \cap C_0(\R)$, by a previous proposition,
    \[
        \widehat{f_n'}(\zeta) = 2 \pi i \zeta \hat{f_n}(\zeta) \text{ for all } n
    \]   
    So by defn,
    \[
        \lim_{n \to \infty} \| \widehat{f_n'} - \widehat{f'} \|_2 = 0 
    \] 
    This implies that $\{\widehat{f_n'}\}$ Cauchy in $L^2(\R)$. Therefore
    \[
        \{2 \pi i (\cdot) \hat{f_n}(\cdot)\}_{n=1}^{\infty} \text{ is also Cauchy in } L^2(\R) 
    \]   
    Hence, $\exists h \in L^2(\R)$ s.t. $2 \pi i \zeta \hat{f_n}(\zeta) \xrightarrow{n \to \infty} h(\zeta)$ in $L^2(\R)$.
    \newline \textbf{Claim:} $h(\zeta) = 2 \pi i \zeta \hat{f}(\zeta)$.
    \newline \noindent We have $2 \pi i \zeta \hat{f_n}(\zeta) \rightharpoonup h(\zeta)$ in $L^2(\R)$.
    So $\forall g \in L^2(\R)$,
    \[
        \int 2 \pi i \zeta \hat{f_n}(\zeta) g(\zeta) \; d\zeta \to \int h(\zeta) g(\zeta) \; d\zeta
    \]   
    We claim that $\forall \phi \in C_c^\infty(\R)$,
    \[
        \int 2 \pi i \zeta \hat{f_n}(\zeta) \phi(\zeta) \; d\zeta \to \int 2 \pi i \zeta \hat{f}(\zeta) \phi(\zeta) \; d\zeta
    \]  
    To prove this,
    \begin{align*}
        \left| \int 2 \pi i \zeta (\hat{f_n}(\zeta) - \hat{f}(\zeta)) \phi(\zeta) \; d\zeta \right| &\leq |2 \pi | \int |\hat{f_n} - \hat{f}| \underbrace{|\zeta \phi(\zeta)|}_{\in C_c^\infty(\R) \Rightarrow L^2(\R)} \; d \zeta  \\
        &\leq |2 \pi | \| \hat{f_n} - \hat{f} \|_2 \underbrace{\| \zeta \phi(\zeta) \|_2}_{= C}  \\
        &= |2 \pi | \| f_n - f \|_2 \| (\cdot) \phi(\cdot) \|_2 \xrightarrow{n \to \infty} 0  
    \end{align*}
    So $\forall g \in L^2(\R)$, $\forall \eta > 0$, $\exists \phi \in C_c^\infty(\R)$ s.t. $\| g - \phi \|_2 < \eta $.
    \begin{align*}
        &\left| \int 2 \pi i \zeta \hat{f_n}(\zeta) g(\zeta) \; d\zeta - \int 2 \pi i \zeta \hat{f}(\zeta) g(\zeta) \; d\zeta \right| 
        \leq \left| \int 2 \pi i \zeta \hat{f_n}(\zeta)(g(\zeta) - \phi(\zeta))  \; d \zeta \right| \\  &+ \left| \int 2 \pi i \zeta (\hat{f_n}(\zeta) - \hat{f}(\zeta)) \phi(\zeta) \; d\zeta \right| + \left| \int 2 \pi i \zeta \hat{f}(\zeta) (\phi(\zeta) - g(\zeta)) \; d\zeta \right| \\
         &\leq \| 2 \pi i (\cdot) \hat{f_n}(\cdot) \|_2 \| g - \phi \|_2 + \| 2 \pi i (\cdot) \hat{f}(\cdot) \|_2 \| g - \phi \|_2 \\
         &\leq 2 C \eta + \lim_{n \to \infty} "" = 0   
    \end{align*}    
    Thus by uniqueness of limits, $h(\zeta) = 2 \pi i \zeta \hat{f}(\zeta)$ which proves $\widehat{f'}(\zeta) = 2 \pi i \zeta \hat{f}(\zeta)$. \qedhere  
\end{proof}

\begin{rem}
    Why is it important to extend the Fourier transform to $L^2(\R)$? 
    Sobolev spaces - a type of weak derivative which belongs to $L^2(\R)$.
    \[
        H'(\R) = \left\{ f \in L^2(\R) \text{ s.t. } f' \text{(weak deriv)} \in L^2(\R) \right\} 
        = \left\{ f \in L^2(\R) : 2 \pi i (\cdot) \hat{f}(\cdot) \in L^2(\R) \right\} 
    \]   
\end{rem}

\subsection{Poisson Summation Formula}
We saw that the Fourier transform is basically some continuous version of Fourier coefficients in
Fourier series. 
All Fourier Transform stuff is consistent with Fourier series. For example, let $f : [0, 1) \to \R$,
and $f \in L^2([0, 1)) \Rightarrow L^1([0, 1))$. Then
\[
    \tilde{f}(x) = f(x) \chi_{[0, 1)}(x) \in L^1(\R)
\] 
Let's periodize $\tilde{f}$ by defining 
\[
    g(x) = \sum_{k \in \Z} \tilde{f}(x - k) = \sum_{k \in \Z} f(x-k) \chi_{[0, 1)}(x-k) = \sum_{k \in \Z} f(x-k) \chi_{[k, k+1)}(x)
\]  
Then
\begin{align*}
    \widehat{\tilde{f}}(k) = \int_{\R} \tilde{f}(x) e^{-2 \pi i k x} \; dx &= \int_0^1 f(x) e^{-2 \pi i k x} \; dx \\
    &= \int_0^1 g(x) e^{-2 \pi i k x} \; dx \\
    &= \hat{f}(k) \quad \text{(Fourier coefficients of $f$)}
\end{align*}
Now our goal is to start with a function $f \in L^1(\R)$ and periodize it and compare to Fourier series.

The first method is shifting:

Consider
\[
    \sum_{k \in \Z} f(x-k) = \sum_{k \in \Z} \tau_k f(x) := Pf(x)
\]  
If the series converges, then
\[
    Pf(x \pm 1) = \sum_{k \in \Z} f( x \pm 1 - k) = \sum_{k \in \Z} f(x - \tilde{k}) = Pf(x)
\] 
Therefore $Pf$ is $\mathbb{T}$ periodic.   

Alternatively, 
\[
    \hat{f}(n) = \int_{\R} f(x) e^{-2 \pi i n x} \; dx
\] 
Then define
\[
    F f(x) = \sum_{n = -\infty}^\infty \hat{f}(n) e^{2 \pi i n x}
\] 
If this series converges, then it is a Fourier series, so it is automaticaly $\mathbb{T}$-periodic.

We will see that the Poisson summation formula relates these two methods of periodization, and thus
relates Fourier series and Fourier transforms.
\begin{thm}
    Let $f \in L^1(\R)$. Then $Pf : \mathbb{T} \to \C$, given by
    \[
        Pf(x) =\footnote{converges pointwise a.e. and in $L^1(\mathbb{T})$} \sum_{k \in \Z} f(x-k)
    \]   
    and $\| Pf \|_1 = \int_0^1 |Pf(x)| \; dx \leq \| f \|_1 = \int_{\R} |f(x)| \; dx$ 
\end{thm}
\begin{proof}
    Let $Q := [-\frac{1}{2}, \frac{1}{2}] \subseteq \R$. Then $\R = \bigcup_{j \in \Z} Q + j$ (disjoint).
    \begin{align*}
        \infty > \int_{\R} |f(x)| \; dx &= \sum_{j \in \Z} \int_{Q + j} |f(x)| \; dx \\
        &= \sum_{j \in \Z} \int_Q \underbrace{|f(x - j)|}_{\geq 0, \text{ Tonnelli }} \; dx \\
        &= \int_Q \sum_{j \in \Z} |f(x-j)| \; dx
    \end{align*}
    $\Rightarrow \sum_{j \in \Z} |f(x-j)| \in L^1(Q) = L^1(\mathbb{T})$ by periodicity. $(*)$ 
    
    \noindent Thus 
    \[
        \int_Q |Pf(x)| \; dx = \int_Q \left| \sum_{j \in \Z} f(x-j) \right| \; dx \leq \int_Q \sum_{j \in \Z} |f(x-j)| \; dx = \| f \|_1
    \] 
    So we have $\| Pf \|_1 \leq \| f \|_1 $.

    By $(*)$ we have for a.e. $x$, $\sum_{j \in \Z} |f(x-j)| < \infty$. So for a.e. $x$,
    \[
        \lim_{N \to \infty} \sum_{|j| > N} |f(x-j)| = 0
    \]  
    Therefore $Pf(x) = \sum_{j \in \Z} f(x-j)$ pointwise a.e.
    
    By $(*)$, 
    \[
        \int_Q \sum_{j \in \Z} |f(x-j)| \; dx = \sum_{j \in \Z} \int_Q |f(x-j)| \; dx < \infty
    \]  
    Therefore
    \begin{align*}
        \lim_{N \to \infty} \sum_{|j| > N} \int_Q |f(x-j)| \; dx &= \lim_{N \to \infty} \int_Q \underbrace{\sum_{|j| > N} |f(x-j)|}_{\leq \sum_{j \in \Z} |f(x-j)| \in L^1(\R)} \; dx \\
        &= \int_Q \underbrace{\lim_{N \to \infty} \sum_{|j| > N} |f(x-j)|}_{= 0} \; dx = 0 \quad \footnote{by dominated convergence theorem}\\
    \end{align*}
    Therefore
    \[
        0 = \lim_{N \to \infty} \int_Q \sum_{|j| > N} |f(x-j)| \; dx \geq \lim_{N \to \infty} \int_Q |Pf - \sum_{j = -N}^N f(x-j)| \; dx
    \] 
    So $Pf = \sum_{j \in \Z} f(x-j)$ in $L^1(Q) = L^1(\mathbb{T})$. \qedhere  
\end{proof}

\begin{lem}
    Let $f \in L^1(\R)$, $Pf \in L^1(\mathbb{T})$. Then $\forall k \in \Z$,
    \[
        \int_{\R} f(x) e^{-2 \pi i k x} \; dx = \hat{f}(k) = \widehat{Pf}(k) = \int_0^1 Pf(x) e^{-2 \pi i k x} \; dx
    \]  
\end{lem}
\begin{proof}
    We check that 
    \[
        \widehat{Pf}(k) = \int_0^1 \sum_{j \in \Z} \underbrace{f(x-j) e^{-2 \pi i k x}}_{\leq |f(x-j)|} \; dx
    \] 
    and $\int_0^1 \sum_{j \in \Z} |f(x-j)| \; dx < \infty$. Therefore by Fubini,
    \[
        \widehat{Pf}(k) = \sum_{j \in \Z} \int_0^1 f(x-j) e^{-2 \pi i k x} \; dx = \sum_{j \in \Z} \int_{-j}^{1-j} f(\tilde{x}) e^{-2 \pi i k (\tilde{x}+j)} \; d\tilde{x}
    \]  
    and $e^{-2 \pi i k j} = cos(-2 \pi k j) + i sin(-2 \pi k j) = 1$. So
    \[
        \widehat{Pf}(x) = \sum_{j \in \Z} \int_{-j}^{1-j} f(\tilde{x})e^{-2 \pi i k \tilde{x}} \; d \tilde{x} = \int_{\R} f(\tilde{x})e^{-2 \pi i k \tilde{x}} \; d \tilde{x} = \hat{f}(k) \qedhere
    \]  
\end{proof}

\begin{thm}[Poisson Summation Formula]
    Let $f \in C(\R)$, and assume $\exists C, \epsilon > 0$ s.t. $|f(x)| \leq C(1 + |x|)^{-(1 + \epsilon)}$. As a consequence of
    this, $f \in L^1(\R)$. Assume $|\hat{f}(\zeta)| \leq C(1 + |\zeta|)^{-(1 + \epsilon)} \Rightarrow \hat{f} \in L^1(\R)$.
    Then $\forall x \in \mathbb{T}$,
    \[
        Pf(x) = \sum_{k \in \Z} f(x-k) = \sum_{k \in \Z} f(x + k) = \sum_{k \in \Z} \hat{f}(k) e^{2 \pi i k x} = Ff(x)
    \]       
    where both series converge absolutely and uniformly on $\mathbb{T}$. 
    In particular, at $x = 0$,
    \[
        \sum_{k \in \Z} f(k) = \sum_{k \in \Z} \hat{f}(k)
    \]  
\end{thm}
\begin{proof}
    Fix $x \in \mathbb{T}$. Check
    \[
        \sum_{|k| > N} |f(x + k)| \leq C \sum_{|k| > N} \frac{1}{(1 + |x + k|)^{1 + \epsilon}}
    \]  
    Now, $|x + k| = |k - (-x)| \geq |k| - |x| \geq |k| - 1$. So
    \begin{align*}
        C \sum_{|k| > N} \frac{1}{(1 + |x + k|)^{1 + \epsilon}} &\leq C \sum_{|k| > N} \frac{1}{|k|^{1 + \epsilon}} \\
        &\leq 2C \int_N^\infty \frac{1}{y^{1 + \epsilon}} \; dy \\
        &= 2C y^{-\epsilon} \Big|_N^\infty = \frac{2C}{N^{\epsilon}} \xrightarrow{N \to \infty} 0
    \end{align*} 
    Therefore
    \[
        \sup_{x \in \mathbb{T}} \left| \sum_{|k| > N} f(x + k) \right| \xrightarrow{N \to \infty} 0
    \] 
    Therefore $\sum_{k \in \Z} f(x + k)$ converges absolutely and uniformly on $\mathbb{T}$.  
    
    Let $P_N f(x) = \sum_{k = -N}^N f(x + k)$. So $P_N f \to P f$ uniformly on $\mathbb{T}$ and $P_N f$ is 
    continuous for all $N$, so $Pf \in C(\mathbb{T}) \Rightarrow Pf \in L^2(\mathbb{T})$. Since $L^2(\mathbb{T})$ is
    a Hilbert space, 
    \[
        Pf \stackrel{\footnote{in $L^2(\mathbb{T})$} }{=}\sum_{k \in \Z} \widehat{Pf}(k) e^{2 \pi i k x}
    \]     
    But by the previous lemma, $\widehat{Pf}(k) = \hat{f}(k)$, so
    \[
        Pf(x) \stackrel{\footnote{in $L^2(\mathbb{T})$}}{=} \sum_{k \in \Z} \hat{f}(k) e^{2 \pi i k x}
    \] 
    By the exact same computation, since $|\hat{f}(\zeta)| \leq \frac{C}{(1 + |\zeta|)^{1 + \epsilon}}$,
    \[
        \underbrace{\sum_{k \in \Z} \hat{f}(k) e^{2 \pi i k x}}_{= F f(x)} \text{ converges absolutely and uniformly on } \mathbb{T} 
    \]  
    So $F_N f(x) = \sum_{k = -N}^N \hat{f}(k) e^{2 \pi i k x}$ and $F_n f \to F f$ uniformly on $\mathbb{T}$. 
    $x \mapsto e^{2 \pi i k x}$ is continuous, so $F_N f$ is continuous for all $N$, and thus $F f \in C(\mathbb{T})$.
    Therefore $F_n f \xrightarrow{L^2(\mathbb{T})} F f$.
    So we have 
    \[
        P_N f = F_N f \xrightarrow{L^2(\mathbb{T})} Pf = F f
    \] 
    Therefore $P f = F f$ a.e.
    
    Finally, since $P f$ and $F f$ are continuous, $P f \equiv F f$ everywhere on $\mathbb{T}$. \qedhere  
\end{proof}

\begin{rem}
    \underline{Application of Poisson Summation:}
    \newline Recall $f(x) = e^{-\pi a x^2} \Rightarrow \hat{f}(\zeta) = e^{- \frac{\pi^2 \zeta^2}{a}}$
    So 
    \[
        \sum_{n \in \Z} e^{\pi a n^2} = \frac{1}{\sqrt{a}} \sum_{n \in \Z} e^{- \frac{\pi^2 n^2}{a}}
    \] 
    If $a$ is small, then the LHS converges slowly, but the RHS converges quickly. So we can use the Poisson summation
    formula to accelerate the rate of convergence of the series. (Useful for computers). 

    In number theory, one considers
    \[
        \mathcal{O}(t) = \sum_{n \in \Z} e^{- \pi n^2 t} \quad \text{ called the (Jacobi) theta function}
    \] 
    By Poisson summation, we get
    \[
        \mathcal{O}(t) = \frac{1}{\sqrt{t}}\mathcal{O}\left(\frac{1}{t}\right)
    \] 
    This property is crucial in the proof of the functional equation of the Riemann zeta function.
\end{rem}

\section{Practice Problems}
Question 1: Let $f \in C_c(\R)$, and suppose $\exists \overline{C} > 0$ and $\alpha \in (0, 1)$ s.t.
$\forall \zeta \in \R$, 
\[
    |\hat{f}(\zeta)| \leq \frac{\overline{C}}{|\zeta|^{1 + \alpha}}
\]     
Prove that $f$ is $\alpha$-Holder continuous, i.e. $\exists L $ s.t.
\[
    |f(x+ h) - f(x) | \leq L |h|^\alpha
\]    
Solution:

Since $f \in C_c(\R)$, $f \in L^1(\R)$ and by the decay on $\hat{f}$, we have $\hat{f} \in L^1(\R)$. Thus by the Fourier inversion formula,
for a.e. $x$,
\[
    f(x) = \int_{\R} \hat{f}(\zeta) e^{2 \pi i \zeta x} \; d\zeta
\]      
But since $f \in C_c(\R)$, the above holds for all $x$. Thus
\begin{align*}
    |f(x+ h) - f(x)| &= \left| \int_{\R} \hat{f}(\zeta) e^{2 \pi i \zeta (x + h)} - \hat{f}(\zeta) e^{2 \pi i \zeta x} \; d\zeta \right| \\
    &= \left| \int_{\R} \hat{f}(\zeta) e^{2 \pi i \zeta x} (e^{2 \pi i \zeta h} - 1) \; d\zeta \right| \quad \footnote{Apply change of variables, $\tilde{\zeta} = \zeta h$ }\\
    &= \left| \int \hat{f}(\frac{\tilde{\zeta}}{h}) e^{2 \pi i \frac{\tilde{\zeta}}{h} x} (e^{2 \pi i \tilde{\zeta}} - 1) \frac{1}{h}\; d\tilde{\zeta} \right|  \\
    &\leq \int_{\R} \left| \hat{f}(\frac{\tilde{\zeta}}{h}) \right| (1 + 1) \frac{1}{|h|} \; d\tilde{\zeta} \\
    &\leq 2 \overline{C} \int \frac{1}{|\frac{\tilde{\zeta}}{h}|^{1 + \alpha}} \frac{1}{|h|} \; d\tilde{\zeta} \\
    &= 2 \overline{C} |h|^{\alpha} \int_{\R} \frac{1}{|\tilde{\zeta}|^{1 + \alpha}} \; d\tilde{\zeta} \leq L |h|^\alpha \qedhere
\end{align*} 
Therefore $f$ is $\alpha$-Holder continuous.


 
